\chapter{어휘 목록}

\section*{일러두기}

\Vocab{A1}{벼}{베를(1-11/5:33) \\
나라글(6-15/43:51) \\
나라글(6-15/45:42) \\
나라기라(6-15/48:44) \\
벼라(6-15/48:44) \\
베배기(6-15/48:44) \\
나락(6-15/48:44) \\
나락(6-15/49:00) \\
나락뼈(1-09/22:27) \\
나락(1-09/22:27) \\
벼(1-09/22:27) \\
나락(1-09/23:57) \\
나라기라(1-09/24:13) \\
나락베ː람(1-09/24:19) \\
나락베ː라(1-09/24:19) \\
나락(1-09/24:19) \\
베농사(2-10/03:35) \\
올ː베싸리라(1-11/45:37) \\
벼ː쌀(1-11/45:37)}{베(3-21/20:03)}
\Vocab{A2}{쌀}{쌀(1-11/17:40) \\
찹쌀(1-11/19:49) \\
올베싸리(1-11/21:49) \\
쌀(6-15/43:51) \\
쌀(1-09/19:19) \\
쌀떼(1-09/22:27) \\
쌀대(1-09/22:27) \\
싸리(1-09/24:50) \\
올베싸리랑거(1-09/07:00) \\
올ː베싸리라(1-11/45:37) \\
쌀(1-11/33:51) \\
찹쌀(1-11/33:51) \\
벼ː쌀(1-11/45:37)}{쌀(5-24/24:26) \\
쌀(3-21/20:40) \\
싸리(3-21/41:57) \\
쌀(4-22/15:56)}
\Vocab{A3}{볍씨}{}{나락(4-22/20:32) \\
나라기라(5-23/20:37) \\
나락(3-21/20:38) \\
나랑으(4-22/20:42) \\
나락(3-21/20:43) \\
나라기고(5-23/20:43) \\
나라기고(4-22/21:08) \\
나락씨를(4-22/22:10) \\
벼씨(3-21/22:16)}
\Vocab{A4}{못자리}{}{}
\Vocab{A5}{보리}{보리(1-11/17:40) \\
보리(1-11/21:49) \\
보리가(1-11/23:13) \\
보리(1-11/23:13) \\
보ː리를(6-15/49:00) \\
보리차(6-15/53:12) \\
보리꼬게(6-15/18:21)}{보리가틍(5-24/23:30) \\
보리는[nɨ̃](5-24/23:36) \\
보릴(5-23/22:41) \\
보리도(4-22/22:43) \\
보리는(4-22/22:53) \\
보리가(3-21/22:55) \\
보리(4-22/15:56)}
\Vocab{A6}{깜부기}{}{}
\Vocab{A7}{조}{}{조ː(4-22/12:43)}
\Vocab{A8}{수수}{}{}
\Vocab{A9}{기장}{}{}
\Vocab{A10}{옥수수}{}{옥쑤수(5-24/22:30) \\
강넹이[e.ĩ](7-25/22:31) \\
깡넹이(5-24/22:32) \\
깡넹이라(5-24/22:32) \\
깡넹이(5-24/22:36) \\
옥쑤수(5-24/22:36) \\
옥쑤수고(5-24/23:16) \\
깡넹이라고(5-24/23:16)}
\Vocab{A11}{새끼(줄)}{사내끼(6-15/43:38) \\
사네끼도(6-15/03:33) }{}
\Vocab{A12}{새참}{셰ː참(1-11/16:33) \\
셰ː차미라고(1-11/16:33) \\
셰꺼리(1-11/16:33) \\
새차믄(1-11/16:45) \\
새ː꺼리(1-11/17:00) \\
새ː차미라는(1-11/17:00) \\
새ː참(1-11/17:00) \\
새ː꺼리가(1-11/17:00) \\
새차미란(1-11/17:00) 새ː꺼리를(1-11/21:00) \\
새참(1-11/21:00)}{}
\Vocab{A13}{절구}{도구통이라(1-11/19:49) \\
절구(1-11/19:49) \\
절구통(1-11/31:53) \\
도구통(1-11/31:53)}{}
\Vocab{A14}{겨}{쌀켜(1-11/19:21) \\
쌀껴(1-11/19:21) \\
밀ː겨(1-11/19:21)}{}
\Vocab{A15}{키}{챙이라(1-11/19:36) \\
챙ː을(1-11/30:10)}{}
\Vocab{A16}{김매기}{}{}
\Vocab{A17}{호미}{}{호미(4-22/23:39)}
\Vocab{A18}{써레}{}{}
\Vocab{A19}{쇠스랑}{}{}
\Vocab{A20}{가래(농기구)}{}{}
\Vocab{A21}{도리깨}{}{}
\Vocab{A22}{멍석}{}{덕썩(7-25/69:28) \\
덕써기(5-24/69:40) \\
덕썩(5-24/69:40)}
\Vocab{A23}{깨}{깨(1-11/29:47)}{들께가룰(5-24/19:10)}
\Vocab{A24}{나무}{나ː무(6-15/51:52) \\
나ː문디(6-15/51:52) \\
나ː무통(6-15/59:41) \\
나무는(1-09/21:13)/) \\
나무가(1-09/21:13) \\
나무(1-11/34:43) \\
대나무(1-11/54:34) \\
나므일ː(6-15/18:04)}{매실나뭐(5-24/7:06) \\
뽕나무[mɯ̥](5-24/7:06) \\
나무까(5-24/12:54) \\
나무가(4-22/31:09) \\
나무여(5-23/31:34) \\
나무(5-23/31:34) \\
나ː무(5-23/14:02) \\
나무(4-22/20:02) \\
벼나무(4-22/20:02) \\
나무(4-22/20:02) \\
나ː무(5-24/71:58) \\
나ː무를(5-24/72:35)}
\Vocab{A25}{나물}{나무ː를(1-11/25:06) \\
나무른(1-11/25:17) \\
녹뚜나무리(1-11/25:27) \\
시금치나물(1-11/25:27) \\
나물로(1-09/08:34) \\
나물보다(1-09/08:34)}{나물도(2/7:47) \\
나물료(2/8:15) \\
나물(2/8:32) \\
나무리여(2/8:38) \\
나무를(5-24/16:08) \\
나물헤서(3-21/32:18) \\
나물(4-22/32:20) \\
나무리(3-21/32:25) \\
나물(3-21/32:25) \\
나무리(4-22/32:34) \\
나물(4-22/20:02) \\
나물헤[ɦ](4-22/17:01) \\
나물헤[ɦ](4-22/17:01) \\
나물헤[ɦ](4-22/17:01) \\
나물또(4-22/17:07) \\
호ː방나물도(3-21/17:09) \\
까ː지나물도(3-21/17:09) \\
호방나물도(4-22/17:15) \\
오이나물또(4-22/17:15) \\
호ː방나물도(4-22/17:20) \\
오이나물도(4-22/17:20) \\
고추나물도(4-22/17:20) \\
나물(4-22/17:42) \\
나무리(3-21/17:50)}
\Vocab{}{취나물}{}{치[tʰʲ]나물(5-24/7:38) \\
치나물(5-24/7:38) \\
치나물(5-24/7:41) \\
치[tʰʲ]나물(5-24/7:47) \\
치나물(5-24/8:32) \\
취나물도(4-22/32:06) \\
취나물도(3-21/32:07)}
\Vocab{}{가지}{}{까ː지(7-25/25:21) \\
가ː지를(7-25/25:29) \\
가지[dʲ]를(5-24/25:30) \\
가지를(5-24/25:30) \\
까ː지다(5-24/25:30) \\
까ː지(5-24/25:30) \\
가지고(5-24/25:30) \\
가지는(5-24/25:36) \\
까ː지나물도(3-21/17:09) \\
가ː지(3-21/17:13) \\
가ː지(3-21/17:13) \\
까지(5-23/17:14) 5-23 사투리는 [까ː지]라고 하심.}
\Vocab{A26}{냉이}{}{냉잉으만(5-24/8:15) \\
냉이르(5-24/8:15) \\
넹이는(5-24/14:14) \\
넹이(5-24/14:14) \\
넹이가(5-24/14:14) \\
넹이라고(5-24/14:30)}
\Vocab{A27}{달래}{}{달레(5-24/8:30) \\
달롱게라게(5-24/14:08) \\
달롱게(5-24/14:08)}
\Vocab{A28}{콩나물}{콩나물보다는(1-11/25:27)}{콩나무리여(5-24/8:38) \\
콩나물도(3-21/17:09)}
\Vocab{A29}{상추}{}{상추(5-24/8:38) \\
상치도(4-22/16:16) \\
상추(3-21/16:58)}
\Vocab{A30}{오이}{}{오이는(7-25/14:36) \\
오이(5-24/14:38) \\
오이도(5-24/14:38) \\
오이라(7-25/14:43) \\
웨ː(7-25/14:43) \\
에ː(5-24/14:47) \\
오 이(5-24/14:51) \\
오이나물또(4-22/17:15) \\
오이나물도(4-22/17:20) \\
오이(4-22:17:28) \\
오이를(4-22:17:32)}
\Vocab{A31}{부추}{}{부추(5-24/8:32) \\
부추는(5-24/15:04) \\
솔ː(5-24/15:04) \\
부추고(5-24/15:04) \\
정구지고(5-24/15:04) \\
정구지(5-24/15:12) \\
부추(5-24/15:12) \\
솔ː(5-24/15:12)}
\Vocab{A32}{진달래}{}{진달래(7-25/8:52) \\
진달레(5-24/8:55) \\
진달래(5-24/8:55)}
\Vocab{A33}{철쭉}{}{철쭈기요(5-24/9:08) \\
철쭉(5-24/9:08) \\
철쭈꼬싱갑따(5-24/9:08) \\
철쭈꼳(5-24/9:08)}
\Vocab{}{꽃}{꼳(1-11/11:29) \\
꼬시(1-11/23:13)}{꼬슨(5-24/8:55) \\
꼬시(5-24/9:00) \\
꼬시(5-24/9:08) \\
꼬슨(5-24/9:08) \\
꼬시더라(5-24/9:08) \\
꼬(4-22/19:36) \\
딸기꼳(4-22/19:36) \\
꼰나무(5-23/31:06) \\
꼬까승(5-23/1:02:34)}
\Vocab{A34}{삘기}{}{삐ː비(7-25/9:37) \\
삐ː비(7-25/11:15)}
\Vocab{A35}{머루}{}{삼포도(7-25/10:26) \\
삼포도(5-24/10:30) \\
머루(5-24/10:30) \\
머뤼(7-25/10:36) \\
머루(5-24/10:37) \\
머ː루는(7-25/15:45) \\
머ː루하고(7-25/15:51) \\
머루가(5-24/15:54) \\
머루다(7-25/15:55)}
\Vocab{}{머위(머윗대)}{}{머굳때(5-24/15:49) \\
머굳떼(5-24/15:54) \\
머구떼하[ɦ]고는(7-25/15:55) \\
머구[ɣ]ː세~는(5-24/16:08)}
\Vocab{A36}{다래}{}{다레라그요(5-24/16:36) \\
다레(7-25/16:36) \\
산따레구나(7-25/16:36) \\
다레가(5-24/16:40) \\
다레(5-24/16:47) \\
산다레(7-25/16:48) \\
산따레(5-24/16:49)}
\Vocab{A37}{고욤}{}{}
\Vocab{A38}{복숭아}{}{복쑹아(5-24/11:56) \\
복썽(5-24/12:07) \\
복썽이라헤[ɦ](5-24/12:07) \\
복쑹아(5-24/12:16) \\
복 성(5-24/12:16) \\
복쑹아(5-23/12:50) \\
복썽하[ɦ](5-23/13:24) \\
복쑹아(5-23/13:24)}
\Vocab{A39}{배(과일)}{배(1-09/08:08) \\
배를(1-09/08:28) \\
배로(1-11/42:30) \\
푸빼라(1-11/42:30)}{베(7-25/12:36) \\
베가[ɣ]끼도(7-25/12:36) \\
베가치(7-25/12:45) \\
배(5-24/12:46) \\
베가(5-24/12:46) \\
베(5-24/12:46)}
\Vocab{}{돌배}{}{돌베(5-24/12:46) \\
돌베가테(5-24/12:46)}
\Vocab{A40}{밤(열매)}{밤ː(1-11/45:37)}{바ː미고(2/12:22)}
\Vocab{A41}{뿌리}{}{뿌리를(5-24/13:05) \\
뿌리를(5-24/13:11) \\
(뿌리가)(5-24/13:28) \\
뿌리(4-22/20:02)}
\Vocab{A42}{무}{무ː(6-15/57:35)}{무ː(5-24/12:22) \\
무시(5-24/12:22) \\
무시(5-24/12:28) \\
무시라(5-24/12:28) \\
무ː라(5-24/12:28)}
\Vocab{}{고추}{고츄도(1-11/36:43)}{꼬ː치라(7-25/25:15) \\
꼬ː치(7-25/25:15) \\
꼬치(5-24/25:15) \\
꼬ː추(7-25/25:21) \\
고추는(5-24/25:36) \\
고추나물도(4-22/17:20) \\
고치장(5-23/17:35) \\
고추(5-23/17:35) \\
고추장에다(5-23/17:35) \\
고치까리는(5-24/51:55)}
\Vocab{}{}{}{}
\Vocab{B1}{부엌}{정ː제(6-15/54:21) \\
부어기라고(6-15/54:21) \\
부어기다고(6-15/54:21) \\
부억뽀다도(6-15/54:21) \\
부어기라(6-15/54:53) \\
부어게다(6-15/55:22) \\
부억(1-09/11:33) \\
주방(1-09/11:33) \\
부어글(1-09/15:28) \\
부억(1-09/15:28)}{정게(4-22/25:03) \\
정제(3-21/25:06) \\
정제라이(5-23/25:27) \\
정제(5-23/25:27) \\
정제가(5-23/25:27) \\
정게여(4-22/25:36) \\
정게고(4-22/25:36) \\
정제아(5-23/25:39)}
\Vocab{B2}{솥}{가마솓테따가(6-15/50:28) \\
소슬(6-15/50:47) \\
가마소시라고(6-15/50:47) \\
꺼ː믕소시(6-15/50:47) \\
소딴지(6-15/51:23) \\
솓(6-15/51:23) \\
소테다가(6-15/51:23) \\
솓(1-09/18:46) \\
소시라(1-09/18:46) \\
소두방이(1-09/19:02) \\
부뚜방(1-09/19:19) \\
소뚜방(1-09/19:19)}{서딴지[ʝ]라(5-24/20:47) \\
서딴지[ʝ](5-24/20:47) \\
서딴지[ʝ]라(5-24/20:50) \\
서딴지(5-24/20:50) \\
소테다(3-21/26:33) \\
솓(3-21/26:33)}
\Vocab{B3}{아궁이}{부사키고(6-15/52:35)}{아궁이(4-22/26:39) \\
부삭(4-22/26:39) \\
부삭(4-22/26:46) \\
부사기여(5-23/26:52) \\
부삭(5-23/26:52)}
\Vocab{B4}{부지깽이}{부지뗑이(1-09/13:46) \\
부지뗑이(1-09/13:46) \\
부지뗑이라…요(1-09/14:05)}{부지뗑이로(4-22/27:02) \\
부지껭이로(4-22/27:02) \\
부지뗑이가(5-23/28:04)}
\Vocab{B5}{부삽}{}{}
\Vocab{B6}{화로}{}{}
\Vocab{B7}{땔감/땔나무}{}{}
\Vocab{B8}{베개}{}{}
\Vocab{B9}{의자}{}{으ː자(5-24/69:01)}
\Vocab{B10}{마루}{}{마루(7-25/68:39) \\
마레(5-24/68:43)}
\Vocab{B11}{벽}{베름빡(1-09/27:40) \\
벼름빡(1-09/27:40) \\
베름빠기란(1-09/28:07) \\
베름빠기람말은(1-09/29:16) \\
산ː벼기나(1-11/66:25)}{}
\Vocab{B12}{이엉}{}{}
\Vocab{B13}{주춧돌}{}{주춛똘(4-22/36:11) \\
주춛똘만(4-22/36:11)}
\Vocab{B14}{서까래}{}{}
\Vocab{B15}{변소}{화장실(1-09/30:06) \\
뒤[y]깐(1-09/30:13) \\
똥깐(1-09/30:13) \\
통시(1-09/30:22) \\
해우소(1-09/30:22) \\
칙깐(1-09/31:00) \\
통시(2-10/30:16)}{}
\Vocab{B16}{우물}{세ː미라든지(1-11/36:58)}{우무리(4-22/36:59) \\
우무란나(3-21/37:03)}
\Vocab{B17}{도랑}{}{}
\Vocab{B18}{민물}{}{밈무리(5-24/55:57) \\
밈ː물꼬동(5-24/61:04)}
\Vocab{B19}{광주리}{}{}
\Vocab{B20}{궤짝}{}{}
\Vocab{B21}{맷돌}{매똘(1-11/31:53) \\
매똘(1-11/32:42) \\
맫ː똘로(6-15/44:55) \\
맫ː또레다(6-15/44:55) \\
맫ː똘도(6-15/44:55)}{}
\Vocab{B22}{체/어레미/도두미}{}{체(7-25/19:45) \\
체(5-24/19:45) \\
체라고레(5-24/19:48) \\
얼기[ɣ]미(5-24/19:54) \\
얼기미(5-24/19:54) \\
얼기미라(5-24/20:02) \\
얼게미(5-24/20:02)}
\Vocab{B23}{뚝배기}{}{투가[ɣ]리(7-25/18:43) \\
투가리(7-25/18:44) \\
툭씨바리(7-25/18:46) \\
툭씨발(7-25/18:49)}
\Vocab{B24}{차곡차곡}{차복차복(6-15/44:12)}{}
\Vocab{B25}{깍두기}{깍ː뚜기(6-15/56:25)}{체ː지(7-25/17:25) \\
체ː[tʰʲ]지(7-25/17:27)}
\Vocab{}{무김치, 동치미}{동ː치미(6-15/57:26) \\
동치미(6-15/57:35) \\
싱근지(6-15/57:35)}{싱곤지[dʲ]만(5-24/17:40) \\
싱곤지(5-24/17:40) \\
싱 곤 지(5-24/17:45) \\
싱근지(3-21/41:57) \\
무ː김치(4-22/15:43)}
\Vocab{B26}{두부}{두부(1-11/31:43) \\
생두부고(1-11/32:42) \\
두부고(1-11/32:42) \\
두부를(6-15/44:55) \\
두부(6-15/44:55) \\
순두부(1-2-2/36:08) \\
순두부는(1-2-2
36:08)}{두부요(5-24/19:26) \\
춘두부까(5-24/19:33) \\
두부(5-24/20:38) \\
순두부허[ɦ]고(5-24/20:38) \\
두부허[ɦ]고(5-24/20:38) \\
두부(5-24/21:28) \\
두부고(5-24/21:28)}
\Vocab{B27}{부침개}{}{부치리(7-25/15:26) \\
부처리(5-24/15:26) \\
부처리라(5-24/15:26) \\
부침게도(7-25/15:33) \\
부칭게라게(5-24/15:34) \\
부체리(5-24/21:50)}
\Vocab{B28}{주걱}{주벅(6-15/53:59) \\
주게라(6-15/53:59)}{주벅[β~w](7-25/20:14) \\
주 벅(5-24/20:15) \\
주거긴데(7-25/20:18) \\
주워기[ɣ](7-25/20:18) \\
주워기라(7-25/20:18) \\
주거기고(5-24/20:21) \\
주벅(5-24/20:21)}
\Vocab{B29}{누룽지}{누룽지(6-15/53:12) \\
누룽밥(6-15/53:35)}{누룽지(3-21/40:10) \\
누룽이(5-23/40:10) \\
누룸밥(4-22/40:10) \\
누룽지(3-21/40:18) \\
누룽지(3-21/40:29) \\
누룸밥(5-23/40:29) \\
누룽지는(5-23/40:39) \\
누룽밥(5-23/40:38) \\
누룸바블(5-23/40:55) \\
깡뎅이(5-23/40:55) \\
깡뎅이라(5-23/40:55) \\
깜빱(4-22/40:55) \\
깡뎅이라(5-23/40:55) \\
깡뎅이(5-23/40:55) \\
깜빠비라(4-22/40:55) \\
깡빠비라(4-22/40:55) \\
깡바비라(4-22/40:55) \\
깜바비라그(3-21/41:07) \\
깜ː밥(3-21/41:07)}
\Vocab{B30}{숭늉}{}{숭님(4-22/40:29) \\
수잉님(4-22/40:39) \\
숭니미라고(4-22/40:46) \\
숭님(4-22/40:50) \\
숭닝(3-21/40:52) \\
숭닝이(4-22/40:52)}
\Vocab{B31}{수제비}{}{수제비(7-25/21:56)}
\Vocab{B32}{튀밥}{}{튀밥(7-25/23:02) \\
티밥(5-24/23:05)}
\Vocab{B33}{식혜/감주}{}{}
\Vocab{B34}{왜간장}{웨간장(1-11/39:34) \\
웨간장이라(1-11/39:34)}{}
\Vocab{B35}{짜다, 쓰다, 맵다, 떫다}{짜지(1-11/40:41) \\
짜다(1-11/41:42) \\
씁쓰렁맏(1-11/42:12) \\
씁쓰러운(1-11/42:12) \\
씁쓸하다(1-11/42:12) \\
쓴걷또(1-11/42:12) \\
쓴ː마시다(1-11/43:07) \\
쓰다(1-11/43:07) \\
쓴마다녀(1-11/43:07)}{짬맏(7-25/24:58) \\
짭짜람(5-24/24:58) \\
짠마시고잉[ĩ](5-24/25:00) \\
메움(5-24/25:20) \\
메운(5-24/25:36) \\
씀ː맏(5-24/25:47) \\
떨ː븐(7-25/26:16) \\
떨ː븜(5-24/26:18) \\
떨븜맏(5-24/26:32) \\
떨ː븐(7-25/26:37)}
\Vocab{}{달다(달콤하다), 시다}{단마시(1-11/42:02) \\
단ː마(1-11/43:07)}{다콤허[ɦ]이(7-25/10:00) \\
담[z]ː맏(5-24/25:47) \\
신ː맏(5-24/25:47) \\
담ː맏(5-24/25:59) \\
심ː맏(5-24/25:59)}
\Vocab{B36}{만들다/만들어라}{만들쟈나(1-11/11:29) \\
만드는(1-11/17:40) \\
만드러서(1-11/25:06) \\
만ː들고(6-15/44:55) \\
맨ː드러서(6-15/58:36) \\
맹글어가꼬(6-15/59:41) \\
만ː드는거슨(1-11/33:39) \\
만드는(1-11/34:43) \\
맹그러서(6-15/10:40)}{만드러농께(5-24/14:30) \\
만든(5-24/21:06) \\
만드러(7-25/21:07) \\
만드냐(5-24/21:11) \\
만드레(5-24/21:16) \\
만든다(5-24/21:16)}
\Vocab{B37}{옮기다}{}{만드러가[ɣ]꼬(5-23/14:02) \\
멘ː드러가꼬(7-25/75:48)}
\Vocab{B38}{고쟁이}{}{}
\Vocab{}{조끼}{}{조키(5-24/29:59) \\
조끼(5-24/29:59) \\
저께(5-24/29:59) \\
저께라고(5-24/29:59) \\
적 게(5-24/29:59) \\
조끼(5-24/30:07)}
\Vocab{B39}{모시}{}{모시옫(7-25/29:28) \\
모시가꼬(5-24/29:36) \\
모시떡(5-24/29:36) \\
모시떠게(5-24/29:36)}
\Vocab{B40}{골무}{}{골미(7-25/27:35)}
\Vocab{B41}{가위}{}{가위[wi](7-25/27:01) \\
가시게(5-24/27:04) \\
가시게요(5-24/27:04) \\
가이고(5-24/27:04)}
\Vocab{B42}{고름(의복)}{}{}
\Vocab{B43}{구멍}{구멍(1-11/57:28) \\
구멍을(1-11/59:46)}{구멍(5-24/30:26) \\
구멍이(5-24/30:26) \\
꾸ː멍(5-24/30:38) \\
꾸ː멍이라(5-24/30:38) \\
꾸멍이(5-24/30:40) \\
빵ː꾸(7-25/30:41) \\
빵ː꾸(5-24/30:44) \\
게ː구녁(7-25/75:38) \\
뻘ː꾸녀글(7-25/75:48)}
\Vocab{}{}{}{}
\Vocab{C1}{가위바위보}{가이바이보(1-11/47:56)}{}
\Vocab{C2}{금(선)/긋다}{}{}
\Vocab{C3}{윷}{윤ː노리가튼(1-11/47:56)}{윧ː(7-25/71:18) \\
윤ː노리(5-24/71:19) \\
윧ː(5-24/71:19) \\
윧ː짝(7-25/71:26) \\
윧ː(7-25/71:26)}
\Vocab{}{도, 개, 걸, 윷, 모}{떼(1-11/49:19) 뛔(1-11/49:19) \\
떼라(1-11/49:19) \\
토라고(1-11/49:19) \\
토(1-11/49:19) \\
개ː(1-11/49:19) \\
걸ː(1-11/49:19) \\
걸ː(1-11/50:36) \\
거른(1-11/50:56) \\
윤ː(1-11/49:19) \\
숟(1-11/49:19) \\
유ː시라기도(1-11/49:19) \\
수시라기도(1-11/49:19) \\
숟(1-11/50:36) \\
숟(1-11/50:39) \\
유ː시라긍께(1-11/51:09) \\
수시라(1-11/51:09) \\
유ː슨(1-11/51:09) \\
모ː(1-11/49:19) \\
모ː(1-11/50:39) \\
모(1-11/50:39)}{떼(7-25/71:41) \\
게ː(7-25/71:41) \\
걸(7-25/71:41) \\
숟쩨(7-25/71:41) \\
수시다(7-25/71:41) \\
수시고(7-25/71:41) \\
모고(7-25/71:41) \\
떼하[ɦ]고(7-25/72:44) \\
게ː[ɦ]하고(7-25/72:44) \\
걸하[ɦ]고(7-25/72:44) \\
숟(7-25/72:44) \\
모(7-25/72:44)}
\Vocab{C4}{썰매}{썰메(1-11/52:39) \\
썰메가(1-11/53:28)}{}
\Vocab{C5}{얼레}{연ː도리라(1-11/53:28) \\
연ː도리(1-11/53:28)}{}
\Vocab{C6}{봄}{보메도(1-11/21:49) \\
보메(1-11/23:13) \\
보메(1-11/23:43) \\
봄(1-11/52:12)}{봄(4-22/41:38) \\
봄(4-22/41:40) \\
봄(4-22/41:43) \\
보메(3-21/43:41)}
\Vocab{C7}{여름}{여르미(1-11/4:11) \\
여르메는(1-11/15:09) \\
여름(1-11/52:12)}{여름(4-22/41:38) \\
여름(4-22/41:40) \\
여름(4-22/41:43) \\
여르미(1-1/5:09)}
\Vocab{C8}{가을}{가을(1-11/52:12)}{가을(4-22/41:38) \\
가을(3-21/41:40) \\
가을(4-22/41:43) \\
가으레(3-21/43:41) \\
가으레는(3-21/43:41) \\
가을(3-21/43:41) \\
가으레(5-23/48:22)}
\Vocab{C9}{겨울}{겨울에가치(1-11/8:01) \\
겨우레(1-11/24:43) \\
겨울(1-11/24:43) \\
겨우레(6-15/57:35) \\
겨울(1-11/52:12)}{겨우레(5-24/14:14) \\
겨을(4-22/41:38) \\
겨울(3-21/41:40) \\
겨을(4-22/41:43)}
\Vocab{C10}{춥다}{추우먼(6-15/55:04)}{춤ː냐(7-25/32:01) \\
춤ː냐(5-24/32:03) \\
춤냐(5-24/32:16)}
\Vocab{C11}{차다, 차갑다}{}{}
\Vocab{}{따듯하다, 따숩다}{따숨물(6-15/59:41) \\
뜨ː따드내서(1-11/34:43) \\
따딋헤야데자나이(1-11/34:43)}{따ː수께(7-25/32:32) \\
따ː수께(7-25/32:36) \\
따따데(5-23/18:51) \\
따따댄(5-23/28:25) \\
따딷(5-23/28:39)}
\Vocab{}{덥다}{}{떰ː냐(7-25/32:01) \\
덤ː냐(5-24/32:03) \\
덤냐(32:16)}
\Vocab{C12}{고드름}{}{고도름(7-25/70:08) \\
고드름(7-25/70:08)}
\Vocab{C13}{눈(날씨)}{}{누ː니(7-25/70:19)}
\Vocab{C14}{소나기}{}{소나기(7-25/70:29)}
\Vocab{C15}{번개, 벼락}{}{벙게(7-25/74:05) \\
벙게(2-6/74:06) \\
벙게지(7-25/74:10) \\
베락또(5-24/74:10) \\
베락(5-24/74:10) \\
벙개가(5-24/74:10) \\
벼라기(74:10)}
\Vocab{C16}{새벽}{새부게(1-11/15:09) \\
세보게(6-15/10:40) \\
세복뿌터(6-15/10:40) \\
세버게(6-15/11:57)}{세벽(7-25/66:09) \\
세벽(5-24/66:11) \\
세벽(5-24/66:30)}
\Vocab{C17}{아침}{아침(1-11/10:31) \\
아침빱(1-11/10:31) \\
아치메(6-15/64:09) \\
아침빠비(6-15/65:05) \\
밤ː낟(6-15/01:00) \\
밤낟ː(6-15/01:44)}{아침(5-24/66:32) \\
아치미냬ː게(7-25/66:45) \\
아치미다(5-24/66:45) \\
아ː침(7-25/67:05)}
\Vocab{C18}{낮}{}{}
\Vocab{C19}{저녁}{저녁(1-11/10:31) \\
저녁빱(1-11/10:31) \\
저녁(1-11/13:56) \\
저ː녁빠비(6-15/65:05)}{저녀게(3-21/44:15) \\
저녀기고(5-24/66:40) \\
저녁(7-25/67:24)}
\Vocab{C20}{밤(夜)}{바메(1-11/23:13) \\
밤차로(6-15/61:46) \\
바ː메(6-15/61:46) \\
밤ː낟(6-15/01:00) \\
밤낟ː(6-15/01:44)}{한밤쭝(7-25/67:36)}
\Vocab{C21}{길다}{길ː고(1-11/14:41) \\
긴(1-11/28:48) \\
깅ː거(1-11/47:56) \\
깅거(1-11/54:34)}{}
\Vocab{C22}{짧다}{}{}
\Vocab{C23}{하루, 이틀, … 열흘}{여를쓱도(6-15/60:36) \\
하로(1-11/62:03)}{하루나(5-24/51:10) \\
하루(5-24/51:55)}
\Vocab{C24}{내일}{낼ː(6-15/64:09) \\
네이리라도(1-11/64:17)}{내일(7-25/64:35) \\
내일(5-24/65:03),내일(5-24/65:09) }
\Vocab{C25}{모레}{}{내일모ː레(7-25/64:40) \\
모ː레(5-24/64:41) \\
내ː일모레(5-24/64:46) \\
모ː레(7-25/64:50) \\
모레(5-24/65:03) \\
모레(5-24/65:09)}
\Vocab{C26}{글피}{}{글페(7-25/64:50) \\
글페(5-24/64:54) \\
고페(7-25/64:56) \\
고페(5-24/64:59) \\
글페(7-25/64:59) \\
고페(5-24/65:03) \\
고페(5-24/65:09) 7-25 선생님께서는 [글페]
5-24 선생님께서는 [고페]를 선호하시는 듯.}
\Vocab{C27}{어제}{}{어저[ʝ]께(5-24/60:40) \\
어저께(5-24/65:16) \\
어저께(5-24/65:40) \\
어저께(5-24/65:40)}
\Vocab{C28}{그저께}{}{그저께[ẽ]가(5-24/60:40) \\
그저ː께(7-25/65:23) \\
그저께(5-24/65:23) \\
그ː저께(5-24/65:40) \\
그저께(5-24/65:40)}
\Vocab{C29}{그끄저께}{}{그ː그제(7-25/65:27) \\
그ː그제께(7-25/65:27) \\
그그저께(5-24/65:30) \\
그그저께(5-24/65:40) [어끄저께]는 며칠 전을 뜻하며 [그그저께]와 구분됨.}
\Vocab{C30}{지렁이}{}{지렁이(7-25/58:51) \\
지렁이가(5-24/58:57) \\
지렝이(5-24/58:57)}
\Vocab{C31}{미꾸라지}{미꾸라지(1-11/6:34) \\
미끄라지라(1-11/6:34)}{미끄라지(5-24/58:57) \\
미끄람지(5-24/59:12) 5-24 사투리는 [미끄람지]라고 하심.}
\Vocab{C32}{개구리}{개구[ɣ]리(1-11/7:52) \\
개구리보다(1-11/8:01) \\
개구리하고(1-11/8:01)}{게구리(7-25/59:20) \\
게구[ɣ]리(5-24/59:22) \\
께구리(5-24/59:22) 5-24 사투리는 [께구리]라고 하심.}
\Vocab{C33}{올챙이}{올챙이(1-11/8:01) \\
올ː챙이(1-11/8:01)}{올체이[ĩ](7-25/59:35) \\
올쳉이(7-25/59:35) \\
올쳉이(5-24/59:36) \\
올쳉이다(7-25/59:44) \\
올쳉이다(7-25/59:47) 5-24
7-25 [복쩽이]는 올챙이의 뜻으로 쓰이지 않음.}
\Vocab{}{복쟁이(복어의 일종)}{}{복쩽이는(5-24/59:54) \\
복쩽이는(7-25/59:58) \\
복쨍이야(7-25/59:58)}
\Vocab{C34}{두꺼비}{뚜꺼비(1-11/8:01) \\
뚜꺼ː비(1-11/8:01) \\
두꺼비(1-11/8:01)}{뚜꺼비[β](7-25/60:16) \\
뚜꺼비(5-24/60:16)}
\Vocab{C35}{거머리}{}{검ː자리(5-24/62:20) \\
검ː자리(5-24/62:20) \\
검ː자리(5-24/62:20) \\
검ː자리(5-24/62:20)}
\Vocab{C36}{달팽이}{}{달ː펭이(7-25/60:46) \\
달펭이(5-24/60:47)}
\Vocab{C37}{다슬기}{}{고동(7-25/60:55) \\
고동(5-24/60:56) \\
고동(5-24/60:56) \\
물끄도[õ](5-24/60:56) \\
밈ː물꼬동이구[ɣ]마(5-24/61:04) \\
고동[ð]이(5-24/61:15) \\
고동도(5-24/61:15)  5-24 [고동]은 민물에 사는 것과 바다에 사는 것을 포괄하는 넓은 범주고
전자만을 가리키는 말로 [밈ː물꼬동]이 있으며 후자만을 가리키는 말로 [바닥 고동]이 있으나 통칭 [고동]을 많이 쓰는 듯.}
\Vocab{C38}{우렁이}{}{우렁(5-24/64:01) \\
우렁(5-24/64:01)}
\Vocab{C39}{새끼(동물)}{세끼(1-1/2:09) \\
셰끼(1-1/2:09) \\
세끼라(1-1/2:09) \\
뒈ː지셰끼(1-1/2:09) \\
게ː세끼(1-1/2:09) \\
염ː소세끼(1-1/2:09) \\
세끼를(1-1/2:36)}{지세끼(5-24/63:44) \\
새끼라고(5-24/63:44)}
\Vocab{C40}{벌레}{}{벌레(5-24/58:57) \\
벌레(5-24/61:25) \\
벌레가(5-24/62:02) \\
벌레가(5-24/64:11)}
\Vocab{C41}{서캐}{쎼까리가(6-15/60:11) \\
쎼까리(6-15/60:11)}{쎄까리(4-22/29:48) \\
쎄까리(3-21/29:50) \\
쎄까리(5-23/29:55) \\
색깔이(3-21/29:58) \\
쎄까리(4-22/30:05) \\
쎼까리(5-23/30:05) \\
쎄까리(4-22/30:05)}
\Vocab{C42}{파리}{}{파리구마(7-25/61:30) \\
파리(5-24/61:31) \\
퍼[ɔ]ː리(7-25/61:54) \\
파리(5-24/61:56) \\
파리구마(7-25/64:11) \\
파리느(5-24/64:11)}
\Vocab{C43}{구더기}{}{}
\Vocab{C44}{벼룩}{}{벼루기고(7-25/63:38)}
\Vocab{C45}{벌}{}{볼[β](7-25/13:36) \\
벌ː(7-25/13:36) \\
벌ː(5-24/13:44) \\
벌ː(5-24/13:47) \\
벌ː(7-25/61:38) \\
벌(5-24/61:38) \\
벌ː(5-24/61:38)}
\Vocab{C46}{쥐}{}{지(5-24/63:44) \\
지세끼(5-24/63:44) \\
지(5-24/63:44) \\
지세끼(5-24/63:44)}
\Vocab{C47}{돼지}{뒈ː지(1-11/0:53) \\
뒈ː지셰끼(1-11/2:09) \\
뒈지(1-11/4:11) \\
데ː지(1-11/17:40) \\
뒈ː지보끄밍가(2-10/01:02) \\
뒈ː지(2-10/01:02)}{데지가(3-21/32:35) \\
산뙈지(3-21/32:44) \\
메뙈지(3-21/32:55) \\
메뙈지(4-22/33:00) \\
메뙈지(3-21/33:02) \\
메뙈지(4-22/33:02) \\
데ː지(3-21/44:15) \\
데지는(3-21/44:45) \\
데지아고(3-21/45:04) \\
데지(3-21/45:08) \\
돼ː지(3-21/45:17) \\
돼ː지고기(5-23/16:23) \\
데지가(5-23/16:34) \\
돼ː지(7-25/74:44) \\
돼ː지(5-24/75:01) \\
돼ː지(5-24/75:01)}
\Vocab{C48}{고양이}{}{고양이(5-24/63:21) \\
고양이지로(5-24/63:21) \\
고양이(5-24/63:29) \\
꼬양이(5-24/63:29) \\
괴[ø̃]ː이(7-25/63:35) \\
굉이(7-25/63:35) \\
괭이(7-25/63:38)  5-24 선생님께서는 [고양이]가 표준어라는 의식을 가지고 계신 듯 보였으나
방언형을 잘 이끌어내지 못하셨음. [꼬양이]는 본 방언에서 흔한 두음 경음화에 이끌린 hyperdialectism으로 보임.}
\Vocab{C49}{개}{개ː(1-1/0:53) \\
중께라(1-11/1:44) \\
게ː세끼(1-1/2:23)}{게ː가(5-24/75:29) \\
강아지가(5-24/75:29) \\
개ː(5-24/62:49)}
\Vocab{C50}{여우}{}{여우(7-25/62:43) \\
여우여(5-24/62:49) \\
여우가[ɣ]쓰요(5-24/62:49) \\
여우다(5-24/62:49) \\
여우(5-24/62:49),}
\Vocab{C51}{꿩}{}{꽁(5-23/33:48) \\
수꽁(4-22/33:56) \\
앙꽁(4-22/33:56) \\
꽁(4-22/34:32) \\
꿩ː(5-24/62:49)}
\Vocab{C52}{매}{}{7-25/63:06 '뭐라 그냐 자ː근 독쑤리'라고 하심. [독쑤리]가 매를 포괄할 가능성도.}
\Vocab{}{독수리}{}{독쑤리(7-25/63:02) \\
독쑤리요(5-24/63:05) \\
독쑤리(7-25/63:06) \\
독쑤리(7-25/63:13)}
\Vocab{C53}{말(동물)}{}{말(5-23/44:59) \\
마르(3-21/45:04) \\
말(5-24/62:49)}
\Vocab{C54}{마구간}{마ː구까니라고(6-15/50:00) \\
마ː구에(6-15/50:11) \\
마ː구(6-15/50:11) \\
마ː구까네서(6-15/50:11)}{마ː구깐(7-25/74:29) \\
마ː 구 깐(5-24/74:30)}
\Vocab{C55}{소}{소는(1-11/2:36) \\
소를(1-1/2:36) \\
소(1-11/4:11) \\
소가(1-11/15:09) \\
소도(1-11/17:40) \\
소(6-15/49:00),}{소(3-21/44:15) \\
소를(3-21/44:15) \\
소넌(3-21/44:45) \\
소(3-21/44:59) \\
소(3-21/45:04) \\
소하고(3-21/45:04) \\
소늠(4-22/45:09) \\
소(4-22/45:12) \\
소(5-23/20:12) \\
소(4-22/20:14) \\
소(4-22/20:18)}
\Vocab{C56}{외양간}{마ː구까니라고(6-15/50:00) \\
마ː구에(6-15/50:11) \\
마ː구(6-15/50:11) \\
마ː구까네서(6-15/50:11)}{마게(4-22/45:09) \\
마구(4-22/45:09) \\
소마구가(4-22/45:09) \\
소마구(4-22/45:09) \\
소ː마구(4-22/45:12) \\
소마구(4-22/45:12) \\
소막(7-25/74:37)}
\Vocab{}{닭}{닥(1-11/0:53)}{닥하고(4-22/44:15) \\
닥(4-22/44:15) \\
다가고(3-21/44:15) \\
세닥(4-22/44:15) \\
세다글(4-22/44:15) \\
다글(3-21/44:15) \\
닥(3-21/44:45) \\
닥(4-22/44:45) \\
닥또(4-22/44:45) \\
다가고(3-21/45:04) \\
다가고(4-22/45:08) \\
닥(4-22/45:17) \\
다근(4-22/45:17) \\
닥또(4-22/45:32) \\
다기(5-23/45:59) \\
닥ː또(5-23/16:23) \\
닥(5-23/16:23)}
\Vocab{}{돼지우리}{}{돼ː지마게서(4-22/45:17)  4번 화자도 '돼지막'을 사용함.
돼ː지 우리(7-25/74:44) \\
돼ː지 마기라고(5-24/75:01) \\
돼ː지막(5-24/75:01) \\
돼ː지마기(5-24/75:15)  7-25 선생님께서는 [돼ː지 우리]
5-24 선생님께서는 [돼ː지막]을 선호하심. 이때 5-24 선생님께서 당신과 7-25 선생님의 나이차를 지적하며 [돼ː지막]이 표준말이고 [돼ː지 우리]가 사투리겠다 말씀하신 점이 주목됨. 5-24 선생님의 다른 답변 기운데에도 7-25 선생님의 언어를 의식한 것이 있을 수 있음에 요주의.}
\Vocab{}{닭장}{}{닥짱을(3-21/44:15) \\
닥짱을라(3-21/44:15) \\
닥짱(3-21/44:15) \\
닥짱에서(4-22/45:17) \\
닥짱(3-21/45:46) \\
닥짱(7-25/74:54)}
\Vocab{}{}{}{}
\Vocab{D1}{바다}{바다에서(1-11/6:44) \\
바다에(1-11/9:00) \\
압빠다(1-11/9:40) \\
바다에서(1-11/56:21) \\
바다에서(1-11/58:30) \\
바다에서(1-11/58:37)}{바다에스(4-22/46:15) \\
바다에는(4-22/46:35) \\
바다지(7-25/53:44) \\
바다(7-25/53:44) \\
바다(5-24/53:49) \\
바다라(5-24/53:57) \\
바다에서(7-25/59:58) \\
바다에는(5-24/61:04) \\
바다에도(5-24/61:15) \\
바다에서(5-24/76:08)  7-25 [바닥]은 쓰지 않는다고 함.}
\Vocab{D2}{물고기}{}{고기를(5-24/57:39)}
\Vocab{D3}{-어(魚)}{장어(1-11/6:34) \\
장어나(1-11/6:44) \\
져ː너(1-11/13:56) \\
져ː너는(1-11/13:56) \\
숭어도(1-11/59:46) \\
숭ː어(1-11/59:46) \\
저ː너도(1-11/59:46) \\
홍ː어(1-11/61:08) \\
홍ː어라든지(1-11/61:08)}{숭어는(5-24/50:07) \\
숭ː어레(7-25/50:24) \\
숭어는(5-24/50:25) \\
미녀는(5-24/50:25) \\
미너는(5-24/50:25) \\
미너(7-25/50:29)}
\Vocab{}{망둑어(문절망둑)}{}{문ː저리는(5-24/49:44) \\
문ː저리라고[ɣ](5-24/49:44) \\
문ː저리(5-24/49:44) \\
망ː둥어여(5-24/49:44) \\
문저리가(5-24/49:44) \\
망ː둥어가(5-24/49:44)}
\Vocab{D4}{어부}{}{}
\Vocab{D5}{그물}{그무를(1-11/56:01) \\
그물가치(1-11/57:11) \\
그물로(1-11/61:04) \\
그물(59:46) \\
그물로(61:04)}{그물로(5-24/56:14) \\
그무리라이(5-24/57:36) \\
그무를(5-24/57:39) \\
뜰망(7-25/57:41)}
\Vocab{D6}{아가미}{}{아가미(7-25/53:09) \\
아가미이서지(5-24/53:10)}
\Vocab{D7}{지느러미}{}{지느레미(7-25/53:25)}
\Vocab{D8}{바닷물}{}{바다물(4-22/47:17) \\
바다무리(4-22/47:17) \\
바담무리(5-24/55:57)}
\Vocab{D9}{파랗다}{파레지도(1-11/61:46)}{}
\Vocab{D10}{밀물}{밈ː무리(1-11/9:40) \\
민ː물(1-11/10:19)}{밀ː물(7-25/54:05) \\
밀ː물(7-25/54:59) \\
밀ː물다(7-25/55:02) \\
밀ː무리다(7-25/55:12) \\
밀ː무리라(5-24/55:16)}
\Vocab{D11}{썰물}{썰물(1-11/9:40) \\
썰무리(1-11/9:40)}{썰ː물(7-25/54:10) \\
썰ː무리고(7-25/55:02)}
\Vocab{D12}{갯벌}{뻘(1-11/9:00) \\
뻐리(1-11/9:00) \\
뻐른(1-11/9:21) \\
뻐리라고만(1-11/9:21) \\
뻘빠다기라(1-11/9:40) \\
뻘빠닥(1-11/9:40) \\
개빠닥(1-11/10:31) \\
개빠다게(1-11/10:31) \\
개빠다기라(1-11/10:31) \\
뻐레서(1-11/12:09)}{뻘ː(7-25/56:57) \\
뻘하[ɦ]고(5-24/56:59) \\
뻐ː레(7-25/75:48) \\
뻘ː(7-25/75:48)}
\Vocab{D13}{통발}{통ː빠른(1-11/58:30) \\
통ː바리라(1-11/58:30) \\
통ː발로는(1-11/59:12) 통ː바른(1-11/59:29)}{망ː(5-24/57:48) \\
망ː이라(5-24/57:48) \\
그물망(5-24/57:48)  7-25 어르신께서는 고민하는 듯하셨으나 떠올리지 못하셨다.}
\Vocab{D14}{새우}{새우(1-11/38:24)}{세우(7-25/58:04) \\
세우허[ɦ]고(5-24/58:05) \\
세우도(5-24/58:05)}
\Vocab{D15}{게}{대ː개ː(1-11/61:08) \\
홍게라든지(1-11/61:08) \\
게가(1-11/61:31) \\
청게라(1-11/61:31) \\
청게(1-11/61:31) \\
청ː게라(1-11/61:46)}{게ː(7-25/58:04) \\
게ː허고(5-24/58:04) \\
돌ː게(5-24/58:04) \\
돌ː게가(5-24/58:04) \\
게는(5-24/58:04) \\
꼭께가(5-24/58:04) \\
게ː(7-25/58:30) \\
꼬ː께도(7-25/58:30) \\
둥둥게도(7-25/58:30) \\
농ː게도(7-25/58:30) \\
게ː가(5-24/75:36) \\
기ː(7-25/75:38) \\
돌ː기(5-24/75:40) \\
기ː(5-24/75:40) \\
기ː라써(5-24/75:40) \\
기ː(5-24/75:40) \\
게ː를(5-24/75:40) \\
기ː라고(5-24/75:40) \\
돌ː게(5-24/76:00) \\
꼬께(5-24/76:00) \\
게ː레(5-24/76:00) \\
게ː자[ʝ]버꺼따(5-24/76:00) \\
기ː(7-25/76:05) \\
기ː(5-24/76:05)  특히 7-25 선생님께 [기ː]가 더 보수적인 어형으로 인식되는 듯.}
\Vocab{D16}{바위}{바위[wi](1-11/12:40)}{}
\Vocab{D17}{널배, 뻘배(비표준어)}{}{뻘ː뻬(7-25/76:54) \\
뻘ː떼ː(5-24/76:56)}
\Vocab{D18}{해루질(비표준어)}{}{}
\Vocab{}{}{}{}
\Vocab{E1}{마을}{마으레서(1-11/4:11) \\
마ː을금고에서(6-15/64:09)}{마으레서(1-1/5:23)}
\Vocab{E2}{외국}{외ː국(1-11/63:27) \\
외구게는(1-11/64:43)}{예ː구기라(5-24/42:04) \\
예ː국(5-24/42:04) \\
에ː국(5-24/42:04)}
\Vocab{E3}{어디}{으ː디(1-11/9:00) \\
어디(1-11/10:31) \\
어디가(1-11/28:48) \\
으이서(6-15/45:16) \\
어ː디든지(6-15/46:59) \\
어드가든(6-15/46:59) \\
으ː디로(6-15/46:59) \\
으디로(6-15/63:05) \\
으ː디꼬(6-15/65:05) \\
으디가(6-15/65:05) \\
으ː디(1-11/62:13) \\
으ː디(1-11/64:17) \\
어ː디(1-11/66:25) \\
으디가이쓰(6-15/00:15) \\
어ː디(6-15/06:18) \\
어ː디(6-15/07:19) \\
어디(6-15/19:35)}{어기(5-24/18:05) \\
어디(5-24/39:01) \\
으디(5-24/41:54)}
\Vocab{E4}{모시다/모셔라}{}{}
\Vocab{E5}{하나, 둘, … 열}{다서시(1-11/5:33) \\
여서시도(1-11/5:33) \\
하나야(1-11/15:09) \\
둘(1-11/15:09) \\
다스쌀(1-11/28:37) \\
하나에(1-11/30:10) \\
한나(6-15/44:12) \\
하나(6-15/45:16),
한나(1-09/28:07) \\
하나(1-09/28:07) \\
하나(1-11/49:19) \\
항캉(1-11/49:19) \\
공채그랑건(6-15/02:54) \\
한나른(6-15/07:45) \\
한나(6-15/10:04) \\
하나는(6-15/10:04) \\
한나가지고(6-15/19:35) 두ː게씩(1-11/49:19) \\
두ː캉(1-11/49:19) \\
세ː캉(1-11/49:19) \\
세ː게가(1-11/49:19) \\
세ː캉(1-11/49:19) \\
세게고(1-11/50:56) \\
세ː시라(6-15/15:42) 네ː게(1-11/48:31) \\
네ː캉(1-11/49:19) \\
네ː게(1-11/49:19) \\
네ː게를(1-11/49:19) \\
사ː게절바께(1-11/52:08) \\
다섣카니(1-11/49:19) \\
다섣께고(1-11/50:56) \\
 일고(6-15/06:18) \\
여덥쌀(6-15/06:18) \\
열ː려덥(6-15/04:33) \\
여ː라홉(6-15/04:33) \\
열쌀(1-11/52:56) \\
여레(6-15/06:18)}{일리삼사를(5-24/48:54) \\
하나둘(7-25/48:55) \\
하나둘센ː넫다서(7-25/48:55) \\
일리삼ː사오로(7-25/48:55) \\
하나둘ː센ː넫ː다서여서일고여더라홈[ɦ]녈(5-24/49:03) \\
일리삼사고(5-24/49:03) \\
일리삼사고ː이(5-24/49:13) \\
하나둘[z]섿(5-24/49:17) \\
일리삼사오륙칠팔구십(5-24/49:17) \\
한나(3-21/56:51) \\
두리만(3-21/56:51) \\
한나썩(5-24/59:22) \\
한나(5-24/75:35)}
\Vocab{E6}{스물, 서른, … 아흔}{스무사리여(6-15/04:33)}{}
\Vocab{E7}{에누리}{}{}
\Vocab{}{깎다}{}{깡까조(5-24/45:44) \\
깡ː까꼬(5-24/45:44) \\
깡까써(5-24/46:04) \\
깡까가꼬(5-24/46:04)}
\Vocab{E8}{되}{}{뒈(7-25/46:20) \\
되(7-25/46:20) \\
데(7-25/46:20) \\
데(5-24/46:22) \\
데(5-24/46:43) \\
대성ː떼가(5-24/46:22) \\
소성ː떼가(5-24/46:22) \\
소성ː떼(7-25/46:29) \\
데승ː떼(5-24/46:32) \\
소 승 떼(5-24/46:32) \\
소성ː떼(5-24/46:32) \\
데성떼는(5-24/46:45) \\
소성떼는(5-24/46:45) \\
떼가(5-24/47:00) \\
떼이(5-24/47:07) \\
떼(5-24/47:07) \\
데로(5-24/47:07) \\
데성ː떼로(5-24/47:13) \\
소성떼(5-24/47:13)}
\Vocab{E9}{말(단위)}{}{말(7-25/47:00) \\
데두함말(5-24/47:00) \\
데두함마레(5-24/47:07) \\
데두함마리고(5-24/47:13) \\
마른(5-24/47:13) \\
소성함말(5-24/47:13) \\
말ː로(7-25/47:43) \\
말(7-25/47:43) \\
다섬 말(5-23/14:48) \\
열 말(5-23/14:48)}
\Vocab{}{섬(단위)}{}{섬(7-25/47:36)}
\Vocab{E10}{배(연산)}{베ː가를(1-11/3:45)}{}
\Vocab{E11}{금(금속)}{}{}
\Vocab{}{}{}{}
\Vocab{F1}{위(방향)}{우에가(1-11/8:01) \\
위[wi](1-11/12:40) \\
우에(1-11/12:40) \\
우예(1-11/13:05) \\
우ː게다가(6-15/44:12) \\
우게(6-15/44:12) \\
무뤼[y]에서(1-11/56:02) \\
위[y]에서(1-11/56:21) \\
우ː에서(1-11/61:08) \\
우게(6-15/08:29) \\
우게빵으로(6-15/12:49)}{소누겨(5-24/41:07) \\
우예(4-22/51:44) \\
우게도(5-23/52:28) \\
우게는(5-23/56:25)}
\Vocab{F2}{아래}{}{아레ː(7-25/48:06) \\
아레(7-25/48:06) \\
알ː로(7-25/48:06) \\
아레로(7-25/48:22) \\
알ː로(7-25/48:22) \\
아레고(5-23/52:28) \\
아레(5-23/56:25)}
\Vocab{F3}{머리카락}{머리카락(6-15/59:25)}{머리카락(4-22/58:23) \\
머리카락(4-22/58:29)}
\Vocab{F4}{비듬}{}{비듬(4-22/58:40)}
\Vocab{F5}{눈(신체)}{누느로(1-11/64:17) \\
누느로(1-11/64:43) \\
누느르만(6-15/01:00) \\
누이가(6-15/04:33)}{눈(4-22/59:01) \\
누니(3-21/03:19)}
\Vocab{F6}{귀}{귀[y]에다가(1-11/24:43)}{귀(5-23/59:07)}
\Vocab{F7}{말(언어)}{마른(1-11/1:44) \\
말(1-11/2:36) \\
말로(1-11/5:33) \\
마리(1-11/12:40) \\
말해[ɦ](1-11/13:05) \\
말허먼(1-11/16:19) \\
마린데(1-11/17:00) \\
마를(1-11/17:00) \\
마ː른(1-11/17:40) \\
말함[ɦ](1-11/23:13) \\
마리야(1-11/25:06) \\
마ː리(1-11/26:19) \\
멈ː마린지(1-11/26:19) \\
말ː(1-11/32:42) \\
마른(6-15/50:00) \\
말ː(6-15/53:59) \\
마리여(6-15/62:48) \\
마ː라믄(6-15/63:05) \\
말ː(6-15/63:05) \\
두ː말(6-15/64:09) \\
마리야(1-11/34:43) \\
말하먼(1-11/41:51) \\
고마리야(1-11/47:08) \\
마른(1-11/56:21) \\
마른(1-11/57:59) . 고마리야(1-11/64:43) \\
말허믄(1-11. 66:25) \\
마른(1-11/67:28) \\
말로(6-15/01:44)}{마레만(7-25/4:53) \\
표줌마리(5-24/12:16) \\
표줌머른(5-24/12:28) \\
표줌말(5-24/14:14) \\
표줌마를(5-24/14:30) \\
표준말(5-24/14:38) \\
서울마른(5-24/15:04) \\
표줌마른쟈(5-24/15:26) \\
말ː마(7-25/18:11) \\
말보다(5-24/19:10) \\
말ː항[ɦ](5-24/19:33) \\
표줌마른(5-24/20:21) \\
마례(7-25/23:14) \\
표줌[ʝ]마른(5-24/23:16) \\
표줌마른(5-24/25:30) \\
말(7-25/26:31) \\
표줌말로(5-24/26:32) \\
표줌말(5-24/28:51) \\
표줌말(5-24/30:26) \\
말ː헤[ɦ]서(7-25/30:41) \\
말로(7-25/33:43) \\
존댄말로(7-25/34:57) \\
마를(5-24/36:04) \\
마리야(5-24/36:04) \\
밤ː마를(5-24/36:40) \\
마리(5-24/36:40) \\
밤ː말헐떼는(5-24/36:40) \\
말헤[ɦ]요(7-25/37:43) \\
밤ː마르란허[ɦ]제(7-25/38:49) \\
밤ː마란허[ɦ]제(5-24/38:50) \\
존댐마를(5-24/38:50) \\
존댐말(5-24/39:01) \\
말로(5-24/39:54) \\
존댐마른(5-24/41:27) \\
말하듬(5-24/44:03) \\
마리여(5-24/44:32) \\
말쓰믈(5-24/45:07) \\
마레면(5-24/47:07) \\
마를(7-25/48:22) \\
마른(5-24/49:17) \\
표줌말(5-24/49:44) \\
말ː허고(5-23/59:55) \\
말(5-23/59:55) \\
말ː(5-23/59:55) \\
말ː(4-22/59:58) \\
말ː(5-23/1:00:03) \\
마ː른(7-25/50:48) \\
말ː하[ɦ]자면(5-24/51:55) \\
말히[ɦ]자믄(5-24/52:36) \\
일ː봄마ː리고(7-25/52:50) \\
일봄마ː려(7-25/52:50) \\
마라잠(7-25/54:18) \\
마ː른(5-24/66:16) \\
말ː할[ɦ](5-24/68:14) \\
말ː로(5-24/69:01) \\
말ː헤[ɦ]자믄(7-25/72:44) \\
말ː(5-24/73:23) \\
마ː레(5-24/73:23) \\
마ː르(5-24/76:35)}
\Vocab{F8}{혀}{}{세ː(3-21,5/59:31) \\
혀ː(3-21,5/59:31) \\
헤빠닥(4-22/59:32) \\
쎄빠다기고(5-23/59:33) \\
혀고(5-23/59:33) \\
혀(5-23/59:33) \\
쎄빠닥(4-22/59:35) \\
쎄빠다기요(5-23/59:45) \\
쎄빠닥(5-23/59:45)}
\Vocab{F9}{딸꾹질}{}{떠꾹찔(4-22/1:00:34) \\
떠꺽질(4-22/1:00:34)}
\Vocab{F10}{언청이}{}{}
\Vocab{F11}{하품}{}{하ː품(3-21/1:00:39) \\
하품(4-22/1:00:39) \\
하 품(5-23/1:00:43) \\
하ː품(4-22/1:00:43)}
\Vocab{F12}{가래}{}{}
\Vocab{F13}{목}{}{목(5-23/1:01:15) \\
모가지고(5-23/1:01:15) \\
목(4-22/1:01:15) \\
모가지(4-22/1:01:15) \\
모기라(3-21/1:01:22) \\
모기라(5-23/1:01:30) \\
모가지라(5-23/1:01:30) \\
모가지라하믄(4-22/1:01:33)}
\Vocab{F14}{뺨}{}{빰ː떼린다(4-22/1:01:52) \\
빠ː미(5-23/1:01:59) \\
빰ː도(5-23/1:01:59) \\
빰ː(5-23/1:01:59)}
\Vocab{F15}{주름살}{}{}
\Vocab{F16}{배(신체)}{}{베가(5-24/43:45) \\
배(4-22,5/1:02:12) \\
배가(4-22/1:03:04) \\
배창시가(5-23/1:03:04) \\
베창시라헤(5-23/1:03:04) \\
베창시(5-23/1:03:04) \\
베가(7-25/59:58) \\
베가(7-25/60:05)}
\Vocab{F17}{무릎}{}{무룹(3-21/1:00:11) \\
무룹(5-23/1:00:11) \\
물파기고(5-23/1:00:11) \\
물팍(5-23/1:00:18) \\
무룹(5-23/1:00:18) \\
무릅(4-22/1:00:20) \\
무르팍(4-22/1:00:20) \\
무르비다(3-21/1:00:22) \\
물팍(5-23/1:00:24) \\
물파기라게(5-23/1:00:24) \\
물파기(5-23/1:00:24)}
\Vocab{F18}{뼈}{}{}
\Vocab{F19}{고름(신체)}{}{}
\Vocab{F20}{부풀다}{}{}
\Vocab{F21}{가쁘다}{}{}
\Vocab{F22}{엄살}{}{}
\Vocab{F23}{귀신}{귀시니(1-11/44:20)}{}
\Vocab{F24}{방귀}{}{}
\Vocab{F25}{구린내}{}{}
\Vocab{F26}{고린내}{}{}
\Vocab{F27}{씻다/씻어라}{씨꼬(6-15/54:21) \\
씩꺼서(6-15/54:21) \\
씨꺼가지고(1-11/36:58)}{}
\Vocab{F28}{굵다}{}{굴ː근거여(5-24/19:57) \\
(굴그머는)(5-24/20:02) \\
굴~근(5-24/57:26) \\
굴ː근(5-24/62:02)}
\Vocab{F29}{가늘다}{가능고(1-11/54:34)}{}
\Vocab{}{}{}{}
\Vocab{G1}{매맞다/매맞아라}{}{}
\Vocab{G2}{잃다/잃어버리다}{}{}
\Vocab{G3}{잊다/잊어버리다}{이저부따(1-11/50:39)}{이저브러찌(7-25/4:09) \\
이저불 …(7-25/57:52)}
\Vocab{G4}{가르치다}{갈쳐본저기(1-11/67:39) \\
갈ː치다(6-15/03:33) \\
갈ː친(6-15/04:33)}{가쳐줘야(5-24/42:25) \\
갈쳐(4-22/03:13) \\
가켜주께(5-23/05:06) \\
깔쳐써(5-23/14:34)}
\Vocab{G5}{가리키다}{}{}
\Vocab{G6}{묻다[問]}{무러바봉께(1-11/47:08)}{무러봄(7-25/4:53) \\
무러본(7-25/18:11)}
\Vocab{}{}{}{}
\Vocab{H1}{모친}{엄ː니도(6-15/64:09) \\
엄ː마한테(6-15/64:09) \\
엄ː마(6-15/04:33) \\
엄마가(6-15/07:19) \\
엄마(6-15/14:07)}{엄마(3-21/50:57) \\
엄마(4-22/51:01) \\
어머니(5-23/51:03) \\
엄마라(5-23/51:03) \\
엄니(4-22/51:05) \\
엄마(4-22/51:05) \\
엄니야(5-23/51:10) \\
엄니(5-23/51:10) \\
엄니라(5-23/51:10) \\
엄니(4-22/51:14) \\
엄마는(5-23/51:15) \\
엄니(5-23/51:15) \\
엄니여(5-23/51:20) \\
엄니(5-23/51:20) \\
엄마는(5-23/51:20) \\
엄니에(4-22/51:33) \\
엄니는(4-22/51:39) \\
엄니에(4-22/51:39) \\
엄니(5-23/51:40) \\
엄마(5-23/52:28)}
\Vocab{H2}{모친의 모친}{할무이(6-15/01:44) \\
할멍(6-15/17:19) \\
할멈(6-15/19:35)}{할모니(3-21/51:33) \\
예할머니(4-22/51:36) \\
웨할머니제(4-22/51:39)}
\Vocab{H3}{모친의 부친}{}{하라오찌(3-21/51:33)}
\Vocab{H4}{모친의 손위 여성형제}{}{이모(4-22/51:33) \\
이모(5-23/51:47) \\
이모(5-23/52:28) \\
이모고(5-23/52:28)}
\Vocab{H5}{모친의 손위 여성형제의 남편}{}{이모부(3-21/52:59) \\
이숙(5-23/53:00) \\
이모부(4-22/53:03)}
\Vocab{H6}{모친의 손위 남성형제}{}{에삼초니고(5-23/52:28) \\
삼초니(5-23/52:28) \\
삼촌(5-23/52:28) \\
외삼초는(5-23/54:44) \\
외삼초니야(5-23/54:44) \\
외삼촌(5-23/54:44)}
\Vocab{H7}{모친의 손위 남성형제의 아내}{}{}
\Vocab{H8}{모친의 손아래 여성형제}{}{}
\Vocab{H9}{모친의 손아래 여성형제의 남편}{}{}
\Vocab{H10}{모친의 손아래 남성형제}{}{에삼초니고(5-23/52:28) \\
삼초니(5-23/52:28) \\
삼촌(5-23/52:28) \\
외삼초는(5-23/54:44) \\
외삼초니야(5-23/54:44) \\
외삼촌(5-23/54:44)}
\Vocab{H11}{모친의 손아래 남성형제의 아내}{}{}
\Vocab{H12}{부친}{즈그압지(6-15/46:59) \\
아버지(6-15/46:59) \\
압ː지(6-15/62:48) \\
아부지가(6-15/01:44) \\
아부지가(6-15/04:33) \\
아부지도(6-15/05:40)}{아빠다(3-21/50:57) \\
아빠(3-21/50:57) \\
아빠(4-22/51:01) \\
아부지(4-22/51:05) \\
아빠를(4-22/51:05) \\
아부지(4-22/51:14) \\
아부지(4-22/51:26) \\
우라부지(5-23/11:10) \\
아부지가(5-23/12:50) \\
아버지가(5-23/13:24)}
\Vocab{H13}{부친의 모친}{할멈보고(6-15/63:05) \\
할무이(6-15/01:44) \\
할멍(6-15/17:19) \\
할멈(6-15/19:35)}{할무니(4-22/51:58)}
\Vocab{H14}{부친의 부친}{하라부지는(1-11/21:49) \\
하라버지(6-15/62:48)}{하라브젬(5-24/41:27) \\
하라부지(4-22/51:58) \\
하라브지[ʝ](5-23/07:17) \\
할아부지[ʝ]하[ã]트로(5-23/13:24)}
\Vocab{H15}{부친의 손위 여성형제}{}{고모(5-23/54:02) \\
고모(5-23/54:07) \\
고모(4-22/54:09)}
\Vocab{H16}{부친의 손위 여성형제의 남편}{}{}
\Vocab{H17}{부친의 손위 남성형제}{}{삼초니라(5-23/52:43) \\
크나ː부지(4-22/52:44) \\
크나부지과(4-22/52:51) \\
크나ː부지(4-22/53:46) \\
크나ː부지(4-22/53:49)}
\Vocab{H18}{부친의 손위 남성형제의 아내}{}{크너무니(4-22/53:14) \\
크너무냐(5-23/53:14) \\
크넘니(5-23/53:19) \\
크넘ː니(4-22/53:24) \\
그너메(외부화자/53:26) \\
크너무니란(4-22/53:30) \\
크넘니(4-22/53:30) \\
크넘니(5-23/53:31)}
\Vocab{H19}{부친의 손아래 여성형제}{}{고모(5-23/54:02) \\
고모(5-23/54:07) \\
고모(4-22/54:09)}
\Vocab{H20}{부친의 손아래 여성형제의 남편}{}{}
\Vocab{H21}{부친의 손아래 남성형제}{}{자그나ː부지(4-22/52:44) \\
자근아부지고(5-23/52:44) \\
자그나부지다(5-23/52:46) \\
자그나부지(5-23/52:49) \\
삼촌(5-23/52:49) \\
자그나부지(4-22/52:51) \\
자그나ː부지(3-21/53:20) \\
자그나ː부지(4-22/53:46) \\
자그나ː부지(4-22/53:49) \\
자그나부지라그러고(5-23/53:51) \\
삼촌(5-23/53:51) \\
삼촌(5-23/53:57) \\
자그나버지(5-23/53:57)}
\Vocab{H22}{부친의 손아래 남성형제의 아내}{}{자그넘마(3-21/53:20) \\
자근엄ː니(4-22/53:24) \\
자그넘ː니(4-22/53:24)}
\Vocab{H23}{손위 여성형제}{우리언니(6-15/01:00) \\
언니(07:45)}{언니(5-24/37:45) \\
언니니까(5-24/37:45) \\
언니(5-24/38:50) \\
언니제(5-23/50:36) \\
언니(5-23/50:36) \\
언니(5-23?/50:38) \\
언니(3-21/50:39) \\
언니는(4-22/51:33) \\
언니(4-22/51:44) \\
언니고(4-22/51:44)}
\Vocab{H24}{손위 여성형제의 남편}{}{형부라(5-24/37:45) \\
형부라고(5-24/37:45) \\
형부(5-24/37:45) \\
형부고(5-24/38:11) \\
형부라(7-25/38:40) \\
형부(5-24/38:43) \\
형부(5-23/54:58)}
\Vocab{H25}{손위 남성형제}{오빠(6-15/05:40) \\
오빠드리(6-15/07:45)}{오빠(4-22/49:47) \\
오빠(3-21/49:49) \\
오빠(5-23?/50:38) \\
오빵께(5-23/56:25)}
\Vocab{H26}{손위 남성형제의 아내}{}{올케(4-22/55:09) \\
올케(5-23/55:11)}
\Vocab{H27}{손아래 여성형제}{동ː생들도(6-15/43:17) \\
동생드른(6-15/01:44) \\
동생드를(6-15/01:44)}{동셍(5-23/55:22) \\
동생(5-23/55:33)}
\Vocab{H28}{손아래 여성형제의 남편}{}{제ː낭(4-22/55:41) \\
제낭(4-22/55:41) \\
지ː낭이라(5-23/55:49) \\
지ː낭(5-23/55:49) \\
지낭이나(4-22/55:52) \\
제낭이나(4-22/55:52) \\
지ː낭(4-22/55:52) \\
제ː낭은(4-22/55:52) \\
제ː낭(4-22/55:52)}
\Vocab{H29}{손아래 남성형제}{동ː생들도(6-15/43:17) \\
동생드른(6-15/01:44) \\
동생드를(6-15/01:44) \\
동셍드리(1-11/39:34) \\
동셍들(6-15/02:54) \\
동생드른(6-15/03:33) \\
동ː생드릉(6-15/07:45)}{동셍은(5-24/38:22) \\
동셍(5-24/38:22) \\
남동셍이(5-24/38:22) \\
동셍(5-23/55:22) \\
동생(5-23/55:33) \\
동셍은(5-23/56:25) \\
남동셍은(5-23/56:25)}
\Vocab{H30}{손아래 남성형제의 아내}{}{동상더기라(5-23/55:22) \\
동상덕(5-23/55:22) \\
동상덕(4-22/55:33)}
\Vocab{H31}{아내}{}{부이네(5-24/37:45) \\
부이(5-24/37:45) \\
부인(5-24/38:22) \\
아네는(5-23/55:06) \\
아네(5-23/55:33) \\
아ː내에(5-23/56:29)}
\Vocab{H32}{아내의 모친}{}{어머니미다(7-25/34:45) \\
장ː모니미라(7-25/34:57) \\
장문(7-25/35:01)}
\Vocab{H33}{아내의 부친}{}{}
\Vocab{H34}{아내의 손위 여성형제}{}{처혱으(7-25/37:20) \\
처혱으(7-25/37:23) \\
처혱으(7-25/37:37) \\
처혱(5-24/37:57) \\
처혱이라(5-23/54:18) \\
처 혱(5-23/54:18) \\
처형(4-22/54:24) \\
처 혱(5-23/54:25)}
\Vocab{H35}{아내의 손위 여성형제의 남편}{}{메제(5-24/37:22) \\
형부(7-25/37:23) \\
메제제(5-24/37:28) \\
메제(5-24/37:28) \\
메제고(5-24/37:28) \\
메제고(5-24/37:45) \\
메제고(5-24/37:57) \\
메제(5-24/37:57) \\
메제라(5-24/37:57) \\
메제고(5-24/38:11)}
\Vocab{H36}{아내의 손위 남성형제}{}{매제(4-22/56:21) \\
매제고(5-23/56:25) \\
메제(5-23/56:25)}
\Vocab{H37}{아내의 손위 남성형제의 아내}{}{}
\Vocab{H38}{아내의 손아래 여성형제}{}{처제고잉[ĩ](5-24/38:22) \\
처제는(5-23/54:18) \\
처제라(5-23/54:18) \\
처제(5-23/54:25)}
\Vocab{H39}{아내의 손아래 여성형제의 남편}{}{}
\Vocab{H40}{아내의 손아래 남성형제}{}{처남(5-24/38:22) \\
처나미고(5-23/56:25) \\
처남(5-23/56:25)}
\Vocab{H41}{아내의 손아래 남성형제의 아내}{}{처남덱(4-22/56:37)}
\Vocab{H42}{남편}{남펜(6-15/19:35)}{남ː편(7-25/37:20) \\
남ː펴늘(7-25/37:23) \\
남ː펴는(7-25/37:37) \\
남편한테는(5-24/37:45) \\
남ː펴늠(5-24/37:57) \\
남펴니니까(5-24/38:50) \\
남펀(5-23/53:00)}
\Vocab{H43}{남편의 모친}{씨엄치(6-15/06:18)}{시엄ː니(4-22/09:03)}
\Vocab{H44}{남편의 부친}{}{시아부지(4-22/09:03)}
\Vocab{H45}{남편의 손위 여성형제}{}{}
\Vocab{H46}{남편의 손위 여성형제의 남편}{}{}
\Vocab{H47}{남편의 손위 남성형제}{}{}
\Vocab{H48}{남편의 손위 남성형제의 아내}{}{}
\Vocab{H49}{남편의 손아래 여성형제}{}{}
\Vocab{H50}{남편의 손아래 여성형제의 남편}{}{}
\Vocab{H51}{남편의 손아래 남성형제}{}{}
\Vocab{H52}{남편의 손아래 남성형제의 아내}{}{}
\Vocab{H53}{남편의 남성형제(시아주버니)}{}{}
\Vocab{H54}{딸의 남편}{}{사위(7-25/34:45) \\
사이(4-22/57:32) \\
사이(4-22/57:59)}
\Vocab{H55}{아들의 아내}{메느리는(6-15/06:18) \\
메ː느리가(6-15/14:07)}{메느리지(4-22/57:25) \\
메느리(4-22/57:25)}
\Vocab{H56}{아재}{}{}
\Vocab{H57}{자식의 자식}{솔녀(1-11/21:49) \\
손자들(1-11/21:49) \\
손ː지도(6-15/46:59) \\
손주(6-15/48:15) \\
손지(6-15/65:05)}{손지드른(3-21/48:43) \\
손지느(3-21/48:43) \\
손지는(3-21/48:43) \\
손주들(3-21/48:43) \\
손자(3-21/50:02) \\
손자가(5-23/58:04)}
\Vocab{}{}{}{}
\Vocab{}{거시기}{거서기라(1-11/2:09) \\
거세기(1-11/17:40) \\
거스기(6-15/56:10) \\
거석(6-15/56:55) \\
거서게다가(6-15/58:59) \\
거서근(6-15/64:09) \\
거서거(1-11/62:38) \\
거시기(6-15/12:49) \\
거슥(6-15/14:07) \\
거서글(6-15/18:21) \\
거스기(6-15/19:35)}{거스기다냐(5-24/8:27) \\
거스기요(5-24/12:32) \\
거서기여(3-21/43:41) \\
거서기에서(3-21/45:25) \\
거시ː가(3-21/43:41)}
\Vocab{}{대충대충 친}{}{의자그로(5-23/45:36)}
\Vocab{}{이제[今]}{인자(1-11/1:44) \\
인제(1-11/5:33) \\
이쟈(1-11/16:39) \\
이제(1-11/17:40) \\
인자(6-15/43:17) \\
인쟈(6-15/43:51) \\
인제(6-15/45:42) \\
인제(2-10/04:18) \\
은자는(2-10/04:18) \\
인자(2-10/04:18) \\
인자는(2-10/04:18) \\
으자(2-10/27:04) \\
인쟈(1-11/33:35) \\
인자(1-11/62:03) \\
인자(6-15/01:00) \\
인저(6-15/01:00)}{인쟈[ʝ](5-24/5:41) \\
인자(5-24/7:13) \\
인자(5-24/14:30) \\
인쟈[ĩʎ.ʎã](5-24/16:08) \\
인쟈[dʲ](5-24/16:08) \\
인쟈(5-24/17:40) \\
인냐(5-24/24:26) \\
인저께(5-24/30:44) \\
제(5-24/36:04) \\
인자(5-24/36:04) \\
이자[ʝ](5-24/42:25) \\
이자(5-24/42:35) \\
자(5-24/44:03) \\
인자(5-24/44:32) \\
야(5-24/49:13) \\
냐(5-24/49:13) \\
자[dʲ](5-24/49:17) \\
이자(4-22/22:13) \\
인자(3-21/22:31) \\
인쟈느(5-23/22:45) \\
이쟈(3-21/23:01) \\
이제(5-23/23:47) \\
인제(5-23/23:47) \\
인자(4-22/24:31) \\
은쟈(4-22/24:31) \\
이제(4-22/24:44) \\
인쟈(4-22/24:44) \\
인쟈(3-21/25:15) \\
인자(5-23/26:21) \\
인자(4-22/27:10) \\
인제(3-21/27:45) \\
이제(5-23/28:04) \\
인쟈(5-23/29:23) \\
이쟈(3-21/35:44) \\
은쟈(5-23/40:21) \\
이제(3-21/43:41) \\
인제(4-22/53:49) \\
인제(5-23/1:01:26) \\
헨저(4-22/02:23) \\
쟈[dʲ](3-21/04:10) \\
은쟈(3-21/04:30) \\
은쟈(5-23/05:14) \\
쨔(5-23/12:50) \\
이쟈(5-23/14:19) \\
쟏(5-23/14:19) \\
야(5-23/19:33) \\
쟈(5-24/50:50) \\
쩨(5-24/51:55) \\
쨔(5-24/51:55) \\
이제(5-24/52:55) \\
인져~(7-25/54:18) \\
인제(5-24/55:57) \\
졔[dʲ](5-24/1:02:20)}
\Vocab{}{메주}{메주(1-11/34:43) \\
메주(1-11/34:43) \\
메주를(1-11/34:43) \\
메주를(1-11/36:58)}{}
\Vocab{}{고약하다}{고약허지(1-11/34:43) \\
과ː야게(1-11/34:43) \\
고약허제(1-11/36:15)}{고야게(5-23/24:12)}
\Vocab{}{감}{꼬ː깜(1-09/05:21) \\
감ː(1-11/41:51) \\
푸까미라(1-11/42:02) \\
푿까미라(1-11/42:02) \\
푸까미라(1-11/42:30) \\
푿까미라도(1-11/44:20)}{감ː(1-제삼자/18:17) \\
감ː(1-제삼자/18:28) \\
당ː감(1-제삼자/18:32) \\
오레감(5-23/18:40) \\
오레가믄(5-23/18:40) \\
가ː미고(5-23/18:40) \\
당감(4-22/18:41) \\
감ː도(1-제삼자/18:43) \\
감ː(1-제삼자/18:43) \\
감ː도(5-23/38:43) \\
꼬깜도(5-23/38:43) \\
고깜도(4-22/38:44) \\
꼬ː까믈(5-23/18:40)}
\Vocab{}{도끼}{도치(1-09/21:57) \\
도치로(1-09/22:04)}{}
\Vocab{}{짚을 피우다}{지블 써러서(6-15/50:28) \\
지블 푹푹 쌀ː마가꼬(6-15/50:47) \\
지블(1-09/22:27) \\
피우는(1-09/22:27)}{}
\Vocab{}{싸락대, 쌀대}{싸락떼라(1-09/22:27) \\
쌀떼라(1-09/22:27) \\
쌀대라(1-09/22:27)}{}
\Vocab{}{벼훑이/탈곡기}{홀태로(6-15/44:12) \\
홀태(6-15/44:12) \\
홀테라(1-09/24:19) \\
홀태(1-09/24:19) \\
홀테라(1-09/24:50)}{}
\Vocab{}{쭉정이}{쭉떼기(1-09/24:50) \\
쭉떼기라(1-09/24:50)}{}
\Vocab{}{김치}{물ː김치(6-15/56:10) \\
열ː무김치(6-15/56:10) \\
베ː추김치(6-15/56:10) \\
김치(6-15/56:55) \\
짐치(2-10/09:15) \\
배ː추짐치(2-10/09:17)}{}
\Vocab{}{소나무}{소나무(1-11/21:49) \\
솔ː껭이 나무(2-10/21:21)}{}
\Vocab{}{고추}{고츄도(1-11/36:43) \\
고츄장(1-11/36:43)}{}
\Vocab{}{가라앉다}{까라앙겨(1-11/36:58)}{}